\SourceLeaf{253}%
they determine the coefficients \(\alpha\mathord{.}\ \beta,\ \gamma\mathord{.}\ \delta\). We have
\[
\begin{aligned}
\alpha&=\BareRoot{\frac{E}{E'}}\mathord{.}\frac{\sin(\omega+\omega'-\psi)}{\sin\omega'},&\quad\beta&=\BareRoot{\frac{G}{E'}}\mathord{.}\frac{\sin(\omega'-\psi)}{\sin\omega'}\\
\gamma&=\BareRoot{\frac{E}{G'}}\mathord{.}\frac{\sin(\psi-\omega)}{\sin\omega'},&\delta&=\BareRoot{\frac{G}{G'}}\mathord{.}\frac{\sin\psi}{\sin\omega'}
\end{aligned}
\]
We must add the equations
\[
\begin{aligned}
\cos\omega&=\frac{F}{\BareRoot{EG}},&\cos\omega'&=\frac{F'}{\BareRoot{E'G'}},\\
\sin\omega&=\BareRoot{\frac{EG-FF}{EG}},&\sin\omega'&=\BareRoot{\frac{E'G'-F'F'}{E'G'}}
\end{aligned}
\]
so that the four equations can also be written as
\[
\begin{aligned}
\alpha\BareRoot{(E'G'-F'F')}&=\BareRoot{EG'}\mathord{.}\sin(\omega+\omega'-\psi)\\
\beta\BareRoot{(E'G'-F'F')}&=\BareRoot{GG'}\mathord{.}\sin(\omega'-\psi)\\
\gamma\BareRoot{(E'G'-F'F')}&=\BareRoot{EE'}\mathord{.}\sin(\psi-\omega)\\
\delta\BareRoot{(E'G'-F'F')}&=\BareRoot{GE'}\mathord{.}\sin\psi
\end{aligned}
\]

Since the substitutions \(\mathrm{d}p'=\alpha\mathrm{d}p+\beta\mathrm{d}q\) and \(\mathrm{d}q'=\gamma\mathrm{d}p+\delta\mathrm{d}q\) must transform the trinomial \(E'\mathrm{d}p'^2+2F'\mathrm{d}p'\mathord{.}\mathrm{d}q'+G'\mathrm{d}q'^2\) into \(E\mathrm{d}p^2+2F\mathrm{d}p\mathord{.}\mathrm{d}q+G\mathrm{d}q^2\), we easily obtain
\[
EG-FF=(E'G'-F'F')(\alpha\delta-\beta\gamma)^2
\]
and since, conversely, the latter trinomial must be transformed back into the former by the substitution
\[
(\alpha\delta-\beta\gamma)\mathrm{d}p=\delta\mathrm{d}p'-\beta\mathrm{d}q',\qquad(\alpha\delta-\beta\gamma)\mathrm{d}q=-\gamma\mathrm{d}p'+\alpha\mathrm{d}q'
\]
we find
\[
\begin{aligned}
E\delta\delta-2F\gamma\delta+G\gamma\gamma&=\frac{EG-FF}{E'G'-F'F'}\mathord{.}E'\\
E\beta\delta-F(\alpha\delta+\beta\gamma)+G\alpha\gamma&=-\frac{EG-FF}{E'G'-F'F'}\mathord{.}F'\\
E\beta\beta-2F\alpha\beta+G\alpha\alpha&=\frac{EG-FF}{E'G'-F'F'}\mathord{.}G'
\end{aligned}
\]
\Article{22}
From the general investigation of the preceding article we turn to an application of very wide scope. While \(p\) and \(q\) are still taken in their most general sense, for \(p',q'\) we take the quantities denoted by \(r,\varphi\) in Article 15. These symbols
